\section{Discussion}\label{sec:discussion}

\paragraph{What the framework makes possible}
After a salient external event, the public argument is usually about whether
the event ``caused'' the surge that followed. The framework replaces that
single question with several that have precise answers under a stated model.
How many incidents in the window would have occurred anyway, how many
responded directly to the event, and how many were set off by other incidents
(\cref{thm:attribution})? How much of that cascade passed through online
streams (\cref{prop:at-channel})? How large is the cascade that one incident
starts in each stream, as measured by the multiplier $\Mmat$? Each answer
comes with an interval and with the assumptions it rests on.

The framework is equally explicit about what the data cannot settle. The
autocorrelation of counts cannot tell contagion from clustered shocks
(\cref{thm:equiv}), and an incomplete event calendar or a rigid background
manufactures contagion (\cref{thm:spurious}). Identification has to come from
outside the counts: from a calendar of events fixed blind to the outcome, from
higher-order information that only the full likelihood exploits, or from
provenance links between incidents (\cref{cor:equiv-id}). The analysis plan
is built around that requirement (\cref{sec:plan-events}). The companion
paper enters at a precise point. Recorded incidents are a thinned version of
all incidents, and thinning attenuates apparent contagion in a known way
(\cref{thm:thinning}). A contagion estimate is therefore only as good as the
coverage estimate behind it, and the uncertainty about coverage sets a floor
under its precision (\cref{cor:th-darkfig}). For those who must act on a
forecast, the alert rule makes the trade-off between false alarms and missed
surges an explicit choice, agreed with the people who bear its costs
(\cref{thm:alert}), and the nowcast keeps reporting delays from hiding a surge
that is already under way (\cref{prop:nowcast}).

\paragraph{Limitations}
We stress six.

\emph{Linearity and the absence of inhibition.} In a linear Hawkes process an
incident can only raise the rate of later incidents. Deterrence, protective
deployment after an alert, arrests and saturation all lower it, and the model
cannot represent them. Where they operate, the fitted excitation is a net
effect, and an alert that prompts successful protection makes itself look
like a false alarm (\cref{sec:ethics-risks}). Nonlinear Hawkes processes,
which allow inhibitory kernels through a link function \citep{bremaud1996},
remove this restriction at the cost of the cluster representation on which
the surge decomposition rests.

\emph{Stationarity and constant parameters.} The theory assumes
$\spr(\Amat)<1$ and parameters that do not change over the study period.
Contagion may itself change, for example after a major event or a shift in
how incidents are publicised, and recording probabilities can move with
attention and with breaks in reporting systems (\cref{sec:data-incidents}).
The change-point monitor of \cref{rem:fc-cusum} can detect such shifts but
does not model them. Estimates close to criticality deserve particular
caution, since a misspecified background pushes the apparent branching ratio
towards one (\cref{rem:th-spurious}).

\emph{Pre-specification of events.} Everything rests on the event calendar.
Events missing from it are absorbed into $\Amat$ (\cref{thm:spurious}). The
salience index and its threshold are choices, and the sensitivity analysis can
vary them but not remove them (\cref{sec:plan-sens}). More fundamentally, the
shock term treats event times as exogenous. Some candidate events, such as
public statements, may themselves respond to a wave of incidents. Such
reverse dependence violates the model, and blind coding of the calendar does
not prevent it.

\emph{Classifier error in online streams.} An online stream is the output of
a classifier applied to a defined corpus. It measures that corpus, not
platform-wide prevalence, and its error rates may change after salient events
precisely when they matter most (\cref{eq:data-classerr}). Independent errors
cannot create online-to-offline Granger causality under the null, but they
dilute it under the alternative, and errors that move with common shocks can
mimic it (\cref{rem:th-granger-caveats}). Every online-to-offline coefficient
is a statement about predictive dependence between measured series, not about
the influence of online speech on offline behaviour.

\emph{Ecological interpretation.} The streams are aggregates over recording
areas and offence categories. A coefficient $\alpha_{kj}$ summarises how the
counts of two streams move together in time. It does not say who influenced
whom, and it says nothing about the people who live in a region. The
mechanisms that could produce self-excitation (imitation, emboldenment,
diffusion through networks, repeat targeting) are hypotheses that the model
does not separate (\cref{sec:th-meaning}). Family trees reconstructed by EM
are probabilistic summaries under the model, not findings about individual
incidents.

\emph{Independent recording.} The thinning results assume that each incident
is recorded independently of the history. If salient events raise the
propensity to report, or if severe incidents are both more productive and
more often recorded, the correction can err in either direction
(\cref{prop:th-hetthin}). The plan treats these departures as sensitivity
analyses, not as solved problems.

\paragraph{Extensions}
Four extensions seem most valuable. First, \emph{spatio-temporal kernels} of
the epidemic-type aftershock (ETAS) kind \citep{ogata1988,mohler2011} would
replace a fixed partition into regions by a kernel in space and time
\citep{reinhart2018}. Their outputs would have to respect the ethics protocol,
which rules out street- or venue-level maps (\cref{sec:ethics-risks}). Second,
\emph{marked and typed incidents}: letting the productivity of an incident
depend on its offence type, severity or target addresses the bias of
\cref{prop:th-hetthin} directly, and with many streams, sparse or low-rank
estimates of $\Amat$ \citep{zhou2013,hansen2015} become essential. Third,
\emph{nonlinear Hawkes processes} with inhibitory kernels
\citep{bremaud1996}, whose Granger-causal graphs remain available
\citep{eichler2017}, would let deterrence and protection enter the model
rather than its residuals. Fourth, and closest to this series, \emph{joint
estimation with the dark-figure model}. Here the recording probability enters
as a plug-in from the companion paper. A joint model would treat each list as
a delayed, time-varying thinning of one latent self-exciting process, so that
linked records from several lists inform both the dark figure and the
dynamics. The companion paper names continuous-time, open-population
estimation as its own leading extension \citep{kornsteinlab2026darkfigure};
the framework developed here supplies the point process it needs.

\paragraph{Why model the dynamics}
A count says how many incidents were recorded. It does not say where they
came from. A surge that is mostly a direct reaction to an event calls for
vigilance in the days after similar events. A surge sustained by contagion
calls for breaking the chain early, because every incident prevented also
prevents the incidents it would have triggered. A surge that travels from
online to offline calls for watching the online streams. These are different
claims about how hostility spreads, and they carry different
responsibilities. The data can support some of these claims, and on some
questions they cannot decide between rival explanations. Stating which is
which, with assumptions in the open and a plan fixed before the data are
seen, is how a measurement programme can take seriously both the people who
live through a surge and the evidence about it.
